% SGA 6, Exposición II: fragmento español alineado con la fuente.
% Alcance: apéndice I, conclusión de la sección 1, líneas 4390--4438 inclusive.
% Fuente francesa: sga6_fr_r1/src/fr.tex.
% SHA-256 de la fuente francesa completa:
% 0D2158B9FE3B466309A92165727A0B191422E05FBAA0DACD5806E8322C54BA1A.
\providecommand{\sisi}[2]{#2}

Denotemos por $Z$ el esquema cuasicompacto y cuasiseparado obtenido pegando dos ejemplares de $X$ a lo largo de $U$, y por $j_1,j_2$ las dos inmersiones abiertas canónicas de $X$ en $Z$, de modo que se tiene un cuadrado cartesiano
\[
\begin{tikzcd}
 U \arrow[r,"i"] \arrow[d,"i"'] & X \arrow[d,"j_1"]\\
 X \arrow[r,"j_2"] & Z .
\end{tikzcd}
\]
Se tiene
\[
 j_2^*R^1j_{1*}(\widetilde I)\simeq R^1i_*(i^*\widetilde I),
\]
de donde
\[
 H^0(X,R^1j_{1*}(\widetilde I))\simeq H^1(U,\widetilde I)\ne0,
\]
y por tanto $R^1j_{1*}(\widetilde I)\ne0$, lo que prueba (0.2).

Por otra parte, según (II 3.4 y 3.5), se tiene
\[
 \mathbb RQ\,\mathbb Rj_{1*}(\widetilde I)\simeq j_{1*}(\widetilde I).
\]
Se deduce de ello una sucesión espectral cuya sucesión exacta de términos de grados bajos proporciona el isomorfismo
\[
 R^1j_{1*}(\widetilde I)\simeq R^2Q(j_{1*}\widetilde I).
\]
Como el haz $j_{1*}\widetilde I$ es inyectivo en $\Qcoh(Z)$, se ha probado la primera parte de (0.3).

Finalmente, sea $E\in\Ob\Qcoh(Z)$. Se tiene el isomorfismo de dualidad
\[
 \mathbb R\Hom_{\Mod}(E,\mathbb Rj_{1*}(\widetilde I))
 \simeq \mathbb R\Hom_{\Mod}(j_1^*E,\widetilde I),
\]
pero, como $X$ es afín, se tiene (II 3.5)
\[
 \mathbb R\Hom_{\Mod}(j_1^*E,\widetilde I)
 \simeq \mathbb R\Hom_{\Qcoh}(j_1^*E,\sisi{I}{\widetilde I})
 \simeq \Hom(j_1^*E,\widetilde I).
\]
Se deduce
\[
 \mathbb R\Hom_{\Mod}(E,\mathbb Rj_{1*}(\widetilde I))
 \simeq \Hom(j_1^*E,\widetilde I),
\]
de donde una sucesión espectral cuya sucesión exacta de términos de grados bajos da el isomorfismo
\[
 \Ext^2_{\Mod}(E,j_{1*}\widetilde I)
 \simeq \Hom(E,R^1j_{1*}\widetilde I).
\]
Por consiguiente, el morfismo identidad de $E=R^1j_{1*}(\widetilde I)$ proporciona un elemento no nulo de $\Ext^2_{\Mod}(E,j_{1*}\widetilde I)$; por tanto, el funtor $\mathrm D^b(\Qcoh(Z))\longrightarrow \mathrm D^b(Z)$ no es plenamente fiel, lo que termina la demostración de (0.3).
